\begin{enumerate}
    \item{\textbf{Calcula} la solución óptima de la relajación lineal de este problema:
    
\begin{equation}
\begin{split}
\mbox{max. } z = 5x_1 + 8x_2\\
s.a.:\\
x_1 + x_2 \leq  6\\
5x_1 + 9x_2 \leq 45\\
x_1,\,\,x_2 \in  \mathcal{Z}^+
\end{split}
\label{modeloGomory}
\end{equation}

    A continuación se indica una solución básica y factible de la que se recomienda partir:


\begin{center}
\begin{tabular}{c|c|cccc|}
 & $z$ & $x_1$ & $x_2$ &  $h_1$ & $h_2$ \\ 
 & -40 & $5/9$ & 0  & 0 & -3/4 \\ 
\hline
$h_1$ & 1  & 4/9 & 0 & 1 & -1/9\\ 
$x_2$ & 5  & 5/9 & 1 & 0 & 1/9\\ 
\hline
\end{tabular}
\end{center}

    
    }

    \item{\textbf{Indica} la expresión del corte de Gomory asociado a la variable que tiene parte fraccional de mayor valor en la solución óptima de la relajación lineal obtenida.}
    \item {Añade el corte propuesto y \textbf{obtén} la solución óptima del problema original, problema \eqref{modeloGomory}}
    \item{¿Cuál es la \textbf{expresión de este corte} en función de $x_1$ y $x_2$?}

\end{enumerate}


